\chapter{Ánh Trăng Lừa Dối Của Bố Mẹ}
\markboth{Ánh Trăng Lừa Dối Của Bố Mẹ}{Ánh Trăng Lừa Dối Của Bố Mẹ}

\section{Ba Mẹ Chỉ Muốn Tốt Cho Con}
\begin{truyen}{Ba Mẹ Chỉ Muốn Tốt Cho Con}{Parents Only Want What Is Best}
\chuhoa{H}{ọc sinh:} ``Ba mẹ em lúc nào cũng bảo làm vậy là tốt cho em. Nhưng nhiều lúc em chỉ thấy ngộp thôi.''
\textit{(Student: ``My parents always say they do things for my own good. But a lot of the time, I just feel suffocated.'')}

\teacherpair{Muốn tốt cho con là thật. Nhưng nhiều tình thương của người lớn đi ra ngoài lại mang hình thức kiểm soát, so sánh, và áp đặt. Ý định tốt không tự động sinh ra cách làm tốt. Có những cha mẹ thương con rất thật nhưng vẫn làm con kiệt vì không biết thương thế nào cho đúng.}{``Wanting the best for a child is real. But much adult love comes out in the form of control, comparison, and pressure. Good intention does not automatically produce good method. Some parents genuinely love their children yet still exhaust them because they do not know how to love wisely.''}
\end{truyen}

\section{Kỳ Vọng Được Bọc Trong Hy Sinh}
\begin{truyen}{Kỳ Vọng Được Bọc Trong Hy Sinh}{Expectations Wrapped in Sacrifice}
\chuhoa{H}{ọc sinh:} ``Ba mẹ em hay nói đã hy sinh nhiều cho em nên em phải ráng cho đáng.''
\textit{(Student: ``My parents keep saying they sacrificed a lot for me, so I must succeed to make it worth it.'')}

\teacherpair{Lời đó không sai hoàn toàn. Nhưng nếu hy sinh trở thành sợi dây siết cổ con cái bằng mặc cảm biết ơn, thì nó không còn đẹp nữa. Cha mẹ có quyền mong con biết điều. Nhưng không nên bắt con sống suốt tuổi trẻ như một dự án trả nợ ân tình.}{``That statement is not entirely wrong. But when sacrifice becomes a rope tightening around children through guilt, it loses its beauty. Parents have the right to hope their children appreciate them. But they should not force them to live their youth as a debt repayment project.''}
\end{truyen}

\section{Yêu Thương Mà Không Biết Lắng Nghe}
\begin{truyen}{Yêu Thương Mà Không Biết Lắng Nghe}{Love Without Listening}
\chuhoa{H}{ọc sinh:} ``Có khi em không cần tiền thêm hay lớp thêm, em chỉ cần ba mẹ chịu nghe.''
\textit{(Student: ``Sometimes I do not need more money or more classes. I just need my parents to listen.'')}

\teacherpair{Nhiều bậc cha mẹ rất chăm lo phần vật chất và tương lai, nhưng lại nghèo thời gian nghe con nói thật. Đến khi đứa trẻ lý sự, người lớn tưởng nó hỗn. Có khi đó chỉ là tiếng nói bị dồn nén lâu ngày đang bật ngược ra thôi.}{``Many parents care deeply about material needs and the future, yet remain poor in the time they give to listening. When the child argues back, adults think it is disrespect. Sometimes it is only a long-suppressed voice finally pushing back.''}
\end{truyen}

\begin{baihoc}
Không ít áp lực trong học đường bắt đầu từ nhà.
\textit{(Many school pressures begin at home.)}

Tình thương thiếu hiểu biết rất dễ trở thành một loại áp lực được tha thứ quá lâu.
\textit{(Love without understanding easily becomes a type of pressure forgiven for too long.)}
\end{baihoc}

\begin{ghinhoanh}
\textbf{Short English to remember:}

Love can pressure.

Good intention needs wise expression.
\end{ghinhoanh}

\begin{tuhocnhanh}
1. Em đang mang áp lực học tập của chính mình, hay còn đang gánh cả nỗi lo của bố mẹ? \textit{(Are you carrying only your own academic pressure, or also your parents' anxiety?)}

2. Nếu là phụ huynh, em sẽ thương con theo cách nào để bớt làm nó ngộp? \textit{(If you were a parent, how would you love a child without suffocating them?)}
\end{tuhocnhanh}

\ngancach
